% Part: sets-functions-relations
% Chapter: relations
% Section: operations

\documentclass[../../../include/open-logic-section]{subfiles}

\begin{document}

\olfileid{sfr}{rel}{ops}
\olsection{సంబంధాలపై క్రియలు}

సంబంధాలను మార్చడం లేదా కలపడం తరచుగా ఉపయోగపడుతుంది.
\olref[sfr][rel][ord]{prop:stricttopartial}లో సంబంధాల
\emph{సమ్మేళనాన్ని} పరిశీలించాం. సంబంధాలను జతల సమితులుగా
పరిగణించి వాటి సమ్మేళనం తీసుకోవడమే అది.
అలాగే \olref[sfr][rel][ord]{prop:partialtostrict}లో సంబంధాల
సాపేక్ష భేదాన్ని పరిశీలించాం. సంబంధాలపై చేయగల మరికొన్ని క్రియలు ఇవి.

\begin{defn}\ollabel{relationoperations} 
$R$, $S$ సంబంధాలు అనుకుందాం; $A$ ఏదైనా సమితి కావచ్చు.

$R$ యొక్క \emph{విలోమం}
$R^{-1} = \Setabs{\tuple{y, x}}{\tuple{x, y} \in R}$.

$R$, $S$ల \emph{సాపేక్ష లబ్ధం}
$(R \mid S) = \{\tuple{x, z} : \exists y(Rxy \land Syz)\}$.

$R$ను $A$కు \emph{పరిమితం చేసిన సంబంధం}
$\funrestrictionto{R}{A}= R \cap A^2$.

$R$ను $A$కు \emph{వర్తింపజేసిన ఫలితం}
$\funimage{R}{A} = \{y : (\exists x \in A)Rxy\}$
\end{defn}

\begin{ex}
$S \subseteq \Int^2$ అనేది $\Int$పై ఉత్తరవర్తి సంబంధం అనుకుందాం.
అంటే $S = \Setabs{\tuple{x, y} \in \Int^2}{x + 1 = y}$;
కనుక $Sxy$ కావడానికి అవసరమైన మరియు సరిపడే షరతు $x + 1 = y$ అవడం.

$S^{-1}$ అనేది $\Int$పై పూర్వవర్తి సంబంధం, అంటే
$\Setabs{\tuple{x,y}\in\Int^2}{x -1 =y}$.

$S\mid S$ అనేది
$ \Setabs{\tuple{x,y}\in\Int^2}{x + 2 =y}$.

$\funrestrictionto{S}{\Nat}$ అనేది $\Nat$పై ఉత్తరవర్తి సంబంధం.

$\funimage{S}{\{1,2,3\}}$ అనేది $\{2, 3, 4\}$.
\end{ex}

\begin{defn}[సంక్రామక సంవృతం] $R \subseteq A^2$ ఒక ద్విస్థానిక సంబంధం అనుకుందాం.
	
$R$ యొక్క \emph{సంక్రామక సంవృతం}
$R^+ = \bigcup_{0 < n \in \Nat} R^n$.
ఇక్కడ $R^1 = R$, $R^{n+1} = R^n \mid R$ అని పునరావృత్తంగా నిర్వచిస్తాం.

$R$ యొక్క \emph{స్వావర్తన సంక్రామక సంవృతం}
$R^* = R^+ \cup \Id{A}$.
\end{defn}

\begin{ex}
ఉత్తరవర్తి సంబంధం $S \subseteq \Int^2$ తీసుకుందాం.
$S^2xy$ కావడానికి అవసరమైన మరియు సరిపడే షరతు $x + 2 = y$;
$S^3xy$కు అటువంటి షరతు $x + 3 = y$; ఇలాగే కొనసాగుతుంది.
కనుక $S^+xy$ కావడానికి అవసరమైన మరియు సరిపడే షరతు
ఏదో $n \geq 1$కు $x + n = y$ అవడం.
మరోలా చెప్పాలంటే, $S^+xy$ కావడానికి అవసరమైన మరియు సరిపడే షరతు $x < y$;
$S^*xy$కు అటువంటి షరతు $x \le y$.
\end{ex}

\begin{prob}
$R$ యొక్క సంక్రామక సంవృతం నిజంగానే సంక్రామక సంబంధమని చూపండి.
\end{prob}

\end{document}
